\documentclass[11pt]{article}
\usepackage[utf8]{inputenc}
\usepackage[T1]{fontenc}
\usepackage{lmodern}
\usepackage[margin=1in]{geometry}
\usepackage{csquotes}
\usepackage[style=authoryear,backend=biber,sorting=nyt,maxbibnames=99,giveninits=true,uniquename=init]{biblatex}
\usepackage[hidelinks]{hyperref}
\addbibresource{references.bib}
\begin{document}
\section*{Soil Macroporosity and Drainage-Induced Structure}

\subsection{Overview}

Subsurface (tile) drainage is not a hydraulically passive intervention: by lowering
and stabilising the water table it reorganises the soil pore system, its aeration
status, and the pathways along which water and solutes move to the drain. The
evidence assembled here establishes several well-supported propositions. First,
rapid delivery of water and surface-applied solutes to tile drains is governed by a
small volume of continuous macropores---principally earthworm burrows, shrinkage
cracks, root channels, and old installation scars---rather than by uniform matrix
flow \parencite{Shipitalo2004,Stamm2002,Villholth1998,Zehe2001,Nielsen2015,Petersen2012}.
Second, these flow-active pathways and the broader pore network vary systematically
both with depth and, more weakly, with lateral position relative to the drain line,
so that the bulk soil directly above a drain is often more conductive and more
biologically active than the inter-drain zone \parencite{Shipitalo2004,Alakukku2010,Nielsen2015,Petersen2012,Montagne2009}.
Third, sustained drainage measurably alters static structural properties---reducing
bulk density, increasing air-filled porosity and gas diffusivity, and improving
saturated hydraulic conductivity---although the magnitude of these changes is modest
and frequently smaller than that of tillage or of the passage of time since
installation \parencite{HUNDAL1976,Nakajima2014,LAL1993,Frankl2023,Caron2025}.
Fourth, mechanical structural interventions coupled to drainage---mole drainage and
subsoiling---create large, effective, but time-limited pore networks that dominate
flow in heavy clays \parencite{MULLER1988,HARRIS1993,PLAMENAC1988,JHA1995,Castanheira2010}.
Finally, over decades drainage accelerates pedogenic processes---eluviation, clay
translocation, and iron/manganese redistribution---leaving a spatially heterogeneous,
partly irreversible imprint on soil morphology \parencite{Montagne2008,Montagne2010,Montagne2013,Mercier2000,YliHalla2009}.
The lines of evidence range from field tracer and dye/smoke experiments through soil
physical sampling and micromorphology to numerical dual-domain modelling, and they
converge on a consistent mechanistic picture in which a structured, connectivity-limited
macropore system is the central control on drained-soil behaviour.

\subsection{Macropores and preferential-flow pathways to drains}

\subsubsection{Biopores as the dominant conduits}

Across contrasting soils and climates, the pathways that connect the soil surface to
buried drains are overwhelmingly biological in origin. In a Finnish clay, dye applied
to cracks moved down to the base of the plough layer and then travelled laterally
until it encountered open earthworm burrows, which acted as the only conduits to tile
depth; at 10 cm above the tile, movement to the drain occurred exclusively through
earthworm burrows, of which 106 dyed burrows m$^{-2}$ were counted, and the base of
the plough layer functioned as a hydraulic barrier \parencite{Shipitalo2004}. Mean
infiltration in the bulk soil above the tile lines was about twice that between the
lines and correlated positively with \textit{Lumbricus terrestris} numbers and
biomass; individual burrow infiltration rates were highly variable (6--1043 mL
min$^{-1}$, mean 358 mL min$^{-1}$ with the plough layer removed) and consistently
higher with the plough layer intact (mean 1080 mL min$^{-1}$)
\parencite{Shipitalo2004}. In a Swiss loam, stained structures were concentrated in
the upper profile, earthworm burrows were the most abundant stained features, and
worm burrows were the only structures present with sizable frequency at the depth of
the tile drains \parencite{Stamm2002}. At a German site, roughly 50 worm burrows
m$^{-2}$ reached continuously from the surface to 90 cm depth, where the gravel
embedding of the drains began, and the abundance of such continuous pathways was
identified as the most important (necessary but not sufficient) precondition for
preferential flow of the sorptive herbicide isoproturon \parencite{Zehe2001}.

\subsubsection{Direct visualisation of surface-to-drain connectivity}

Smoke-injection and dye studies provide direct, quantitative confirmation that
continuous air-filled macropores link drains to the surface, and that these openings
are spatially clustered above the drain line. In Denmark, smoke plumes emerged from
earthworm burrows at the surface, and all but one smoke-emitting macropore (SEM) lay
within a 1-m band centred over the drain; strong SEMs marked entrances to more
efficient pathways (4 of 7 strong SEMs traced to the drain versus 1 of 6 weak SEMs,
and 5 of 13 dye applications reached the tile line), and strong SEMs also showed the
highest water-infiltration rates \parencite{Nielsen2015}. A companion study found
smoke fans from open macropores on all observation dates---evidence of continuous
macroporous pathways from the drains to the surface---with the number of SEMs
increasing markedly as the water table dropped from about drain depth to 20--50 cm
below it; reproducible SEM distributions were obtained at air pressures of 0.38--0.93
kPa, most SEMs occurred within a 1.0-m-wide band centred above the drains, and their
distribution was significantly non-random along the drain sections
\parencite{Petersen2012}.

\subsubsection{Quantifying active macroporosity and its size structure}

The pore volume that is hydraulically active is small. In a Danish glacial-till soil,
field analysis of a chloride breakthrough indicated an active macroporosity of only
about 0.17\% of total soil volume, or roughly 0.5\% of total porosity, yet this
minute fraction produced abrupt drainage-concentration increases by factors of 20 and
400 within hours of surface application \parencite{Villholth1998}. Image analysis of
macropore area fraction (MAF) in a Canadian soil showed maximum depth-averaged total
MAF of only 0.0036--0.0087, with the minimum (\textless 0.001) occurring at tile
depth; macropores 0.5--5 mm were most abundant near the surface, 5--10 mm pores
peaked near the top of the B horizon, and where macropores existed at tile depth they
were mostly 0.5--5 mm \parencite{Frey2011}. Fractures contributed only a small share
of total macroporosity (maximum roughly 15\%, 2\% and 7\% at three plots), and
dye staining revealed that \textgreater 20-year-old tile-installation scars were
still channelling infiltrate to the drains, while at one plot macropores appeared to
terminate at a textural boundary above the tile \parencite{Frey2011}. Surface soil
and B-horizon saturated conductivities correlated strongly with MAF ($r = 0.79$ and
$r \geq 0.83$), tying the visible pore network to measured flow capacity
\parencite{Frey2011}.

\subsection{Spatial organisation of pores: depth and distance from the drain}

Two organising gradients recur. Macroporosity and biopore abundance decline sharply
with depth, whereas lateral variation relative to the drain line is real but
generally weaker. In a Finnish clay, macroporosity decreased strongly with depth but
did not differ substantially with distance from the drain; the number and area of
earthworm burrows likewise fell with depth, and root-channel numbers showed a depth
effect with little distance effect \parencite{Alakukku2010}. Nonetheless a lateral
signature was detectable: earthworm burrows were 44\% more numerous above a drain
than at 6--7 m away (an 89 burrows m$^{-2}$ difference, $P = 0.04$), and field
saturated hydraulic conductivity was 31\% lower at 2 m from the drain than on or at
6--7 m from it, while laboratory $K_{\mathrm{sat}}$ declined rapidly with depth but
did not vary clearly with drain position \parencite{Alakukku2010}. At that site slope
explained more variation than distance to the drain, as soil properties and earthworm
abundance increased downslope \parencite{Alakukku2010}. A Canadian study similarly
found no significant effect of tile position on MAF distribution, with the caveat
that one plot may have been biased by wheel traffic \parencite{Frey2011}. Reviews of
drained soils nonetheless report higher infiltration rates and lower bulk density
above drain lines than at mid-drain positions, together with a runoff reduction of
34--55\% relative to comparable undrained soils \parencite{Montagne2009}. The
smoke-injection evidence that flow-active openings cluster within a roughly 1-m band
above drains reinforces the view that drainage induces lateral heterogeneity in
near-surface connectivity even when the static pore volume is laterally uniform
\parencite{Nielsen2015,Petersen2012}.

\subsection{Static structure: bulk density, porosity, aeration, and gas diffusion}

Sustained artificial drainage consistently lowers bulk density and raises air-filled
porosity and gaseous diffusivity, improving aeration. In Ohio, subsurface drainage
produced greater hydraulic conductivity, lower unconfined compressive strength, and
lower surface-crust resistance than undrained treatments, and it decreased bulk
density and increased the volume of air-filled pores at 0.02--1.0 bar suction,
although these pore-volume effects were of smaller magnitude \parencite{HUNDAL1976}.
Drainage significantly reduced bulk density in the 0--10 and 10--20 cm layers (with
total porosity showing the opposite pattern), raised relative gas diffusion
($D_p/D_0$) and air content---especially at 0--10 cm and more so under no-till---and
these gas-diffusion gains were positively correlated with corn grain and stover yield
\parencite{Nakajima2014}. Consistent with an aeration mechanism, artificial drainage
in Georgia increased soil oxygen concentrations and lowered soil moisture, raising net
mineralisation and, more strongly, net nitrification, while denitrification potential
was greater in the wetter undrained plots \parencite{GROFFMAN1984}. Across a large
Canadian survey, low relative gas diffusivity was pervasive (over 93\% of corn fields
below the 0.03 threshold at both 15 and 30 cm), and yield correlated positively with
log$_{10}$ relative gas diffusivity and with saturated hydraulic conductivity at 15
cm; boundary-line analysis suggested maximum yield at the highest measured $D_s/D_o$
(0.16) and a major yield drop below $D_s/D_o$ of 0.0001, underscoring that dynamic
aeration indicators outperform static bulk density in explaining productivity
\parencite{Caron2025}.

Comparisons with undrained or subirrigated soils, however, show the direction of the
bulk-density effect is context-dependent. In one multi-site experiment undrained
plots had lower bulk density, higher aggregate mean weight diameter (MWD), lower
penetration resistance, and greater moisture retention than drained plots, with
overall mean bulk density essentially identical (1.44 versus 1.43 Mg m$^{-3}$) and no
significant drainage effect; the authors attributed lower undrained yields to
transient anaerobiosis rather than to permanent physical damage
\parencite{LAL1993}. In a long-term water-table-management experiment, subirrigated
soils had \textit{lower} bulk density and \textit{greater} total porosity than
continuously drained soils but also lower aggregate stability, smaller MWD and
geometric mean diameter, fewer macroaggregates, and about 11\% greater available
water capacity, while the drainage treatment showed greater penetration resistance,
notably at 30--45 cm \parencite{Baker2004}. Drainage-induced compaction of the surface
was also implicated in reduced organic-matter status: an undrained area retained
higher soil organic matter than the drained area, with drainage reducing organic
matter in the top layers and lowering soil pH overall \parencite{FAUSEY1992}, and
non-drainage increased soil-organic-carbon stock relative to tile drainage by
18--27\% across depths in a related Ohio study \parencite{Kumar2014}. The upper bound
on drainage-based aeration control is set by compaction: in the most compacted soils
all pores remain saturated up to $P = -100$ cm, so aeration control by drainage
becomes impossible and excess-water removal depends almost entirely on evaporative
drying, while plough layers impede downward flow, build perched water tables, and
prevent installed systems from exercising their beneficial effect
\parencite{SMEDEMA1993}.

\subsection{Infiltration, hydraulic conductivity, and structural evolution over time}

Drainage and associated loosening raise infiltration capacity and saturated hydraulic
conductivity, and---importantly---these gains often \emph{develop over years} as
biological and structural processes act on the drier profile. Moling in a Thai acid
sulphate clay increased the average basic infiltration rate from 3.9 $\pm$ 0.3 to
10.5 $\pm$ 1.4 mm h$^{-1}$, decreased bulk density and increased porosity (most
markedly at 25 and 40 cm), and raised air permeability and saturated hydraulic
conductivity, particularly below the plough layer \parencite{JHA1995}. A drainage-age
comparison in Minnesota showed that fields drained more than 15 years earlier had
saturated hydraulic conductivity 7.58 cm h$^{-1}$ faster than fields drained within 5
years ($P < 0.05$), with no difference in bulk density or aggregate stability but
higher biological indicators in the older fields; the authors proposed
preferential-flow development as the likely driver of the conductivity increase and
concluded that soil physical and biological properties evolve over time after
installation \parencite{Frankl2023}. Long-term Iranian paddy trials likewise found
that saturated hydraulic conductivity increased over four years in all layers (largest
in the 0--30 cm layer), with drain discharge and water-table depth increasing over
time, and that shallower drainage improved conductivity more than deeper drainage
\parencite{JafariTalukolaee2016,Talukolaee2018}. Drainable porosity varied seasonally,
being lower during rice than canola seasons \parencite{Talukolaee2018}. Tillage
interacts strongly with these hydraulic changes: a column study showed that after
tillage the tilled column drained earlier (by 11 h) and discharged more (55.2 versus
38.8 mm) than an untilled column, a difference reproduced by HYDRUS only when
tillage effects on hydraulic properties were included \parencite{Moroizumi2004}, and
shallow-tilled sandy soil yielded about twice the transient drainage of deep-tilled
soil, with correspondingly larger solute losses \parencite{MANSELL1980}. Amending the
trench backfill can also govern infiltration: lime-mixed backfill gave a mean
post-steady-flow infiltration rate of 3.2 m d$^{-1}$ versus 2.2 m d$^{-1}$ in control
backfill and 0.83 m d$^{-1}$ between drains, stabilising clay structure and reducing
the strong decline of the drainage coefficient with rising water content
\parencite{Saulys2006}.

\subsection{Mole drainage and subsoiling: engineered structure and its durability}

Where matrix permeability is inadequate, mechanical fissuring by mole drainage and
subsoiling creates the effective pore network. On UK heavy clays, flow through the
system was dominated by mole drains, with up to 90\% of flow occurring through them;
effective drainage deepened the winter water table (drained plots averaged 216 mm
deeper than undrained over 1980--88) and reduced flood risk by 16\% for a 10-year
event, while direct drilling improved subsoil structure enough to extend mole-drainage
life by at least 25\% but also increased surface compaction, runoff, and peak flow
\parencite{HARRIS1993}. On a very deep heavy clay in the former Yugoslavia, subsoiling
and mole drainage raised the soil drainage porosity by up to 20\% in all treated
variants, lowered the critical-period water table to 50 cm within 3 days at 20--30 m
spacing, and---together with permeable trench backfill---eliminated waterlogging
within 2 days; gravel-backfilled trenches retained favourable hydraulic conductivity
and bulk density after 5 years whereas unbackfilled trenches compacted quickly
\parencite{PLAMENAC1988}. The durability of moled structure is finite but
substantial: German long-term trials found mole channels formed above the water table
with an effective top depth $\geq 0.55$ m persisted as open cavities for 3--7 years
and thereafter as zones of increased horizontal permeability, with improved
permeability in disruption zones continuing for more than 54 months, although after 20
years backfilled soil tended to be less permeable than adjacent undisturbed soil
\parencite{MULLER1988}. Combined pipe-plus-mole systems modelled in Portugal drained
1333 m$^3$ ha$^{-1}$ (133 mm) with mole drains versus 1029 m$^3$ ha$^{-1}$ without---a
22\% loss of drainage capacity when moles were absent---and lowered the peak water
table by about 0.1 m \parencite{Castanheira2010}. Mole drainage in Australia increased
gas-filled porosity above the channels (though still below the optimum for roots),
while gypsum-enriched slots further raised slot porosity and improved mole-channel
stability under increased flow \parencite{Muirhead1996}. Deep ripping to 0.45 m with
gypsum, with or without drainage, lowered subsoil bulk density and strength, produced
smaller aggregates, and raised subsoil hydraulic conductivity and air-filled porosity,
with improvements persisting through the two-year study \parencite{Greenwood2006}.
Rice-husk drainage mole systems in Korean paddy soils similarly lowered soil moisture
and raised permeability and aeration, the 2-m interval giving the greatest improvement
and the most uniform moisture distribution \parencite{Jung2024}.

\subsection{Aggregate stability, soil strength, crusting, and trafficability}

Drainage modifies the mechanical behaviour of the surface and subsoil, generally
reducing crusting and improving trafficability while producing mixed effects on
aggregation. Undrained and ponded plots developed strong crusting and adverse
structure, and the undrained treatment had the highest crust density and resistance
and the largest, thickest crust units; crust thickness increased in the order combined
tile-and-surface \textless\ tile \textless\ surface \textless\ undrained
\parencite{HUNDAL1976,FAUSEY1969}. Surface-drained plots required greater penetration
force than plots with subsurface or combined drainage, and combined drainage produced
deeper, more uniformly distributed soybean roots and a significantly longer taproot
\parencite{FAUSEY1969}. Soil strength is tightly coupled to matric suction: in an
English clay, strength at 7.5--22.5 cm depth was linearly related to matric suction
over $-3$ to $-20$ cm water, and a 10-cm rise in suction was on average as effective
at reducing rutting as a 670-kg wheel-load reduction; in ploughed soil, subsurface
drainage reduced tractor-traffic damage, decreasing wheel sinkage and compaction, and
the authors recommended draining clay soils to below 50--60 cm to minimise structural
damage \parencite{STEINHARDT1974}. Trafficability effects were quantified directly in
sodic North Dakota soils: tile drainage lowered drawbar power relative to no drainage
when no amendments were applied, and drainage combined with the lower gypsum rate
produced penetration-resistance results similar to the higher gypsum rate, showing
that drainage can substitute for part of the amendment requirement
\parencite{Wamono2016,Wamono2016b}. Aggregation responses depend on the comparison:
the drainage treatment showed larger MWD and more macroaggregates than subirrigation
\parencite{Baker2004}, no-till plus drainage increased aggregate stability and MWD at
the surface \parencite{Kumar2014}, yet drainage reduced MWD and infiltration relative
to non-drainage in an Iranian trial \parencite{Shabanpour2022}, and drainage-water
recycling raised penetration resistance (1292 kPa) above drainage-only (1197 kPa) and
non-drained (1084 kPa) controls while lowering organic matter and altering texture at
depth \parencite{Kaur2023}. Elevated soil temperature is a further structural benefit:
subsurface drainage significantly warmed soils, with the largest May--July increases
at 30--60 cm and larger spring warming under narrower spacing
\parencite{Jin2008}, consistent with the reported 2--3$^{\circ}$C spring warming and
roughly 10-day-earlier sowing achieved by subsoiling/moling with drainage
\parencite{PLAMENAC1988}.

\subsection{Drainage-induced pedogenesis and long-term structural imprint}

Over years to decades, the redirected water fluxes of a drained profile accelerate
pedological processes and leave a lasting, spatially structured imprint. In a French
Albeluvisol studied by space-for-time substitution, subsurface drainage intensified
eluviation in the E\&Bt and upper Bt horizons, producing lower contents of the
\textless 2 $\mu$m fraction and Fe and higher quartz near the drain, with changes at
60 cm from the drain of the same nature as long-term differentiation but stronger than
those at 110 cm \parencite{Montagne2010}. Drainage removed swelling clay minerals and
more than one-third of the Fe oxides near the drain, increased precipitation of Mn
oxides and Mn-rich ferrihydrite with decreasing distance to the drain, and shifted the
dominant Mn form from Mn-rich ferrihydrite to Mn oxides after 16 years, with
eluviation most active shortly after installation \parencite{Montagne2008}. The
spatial pattern is process-specific and asymmetric: eluviation was largely restricted
to the first metre either side of the drain, whereas redox-driven Mn redistribution
and black-soil-volume formation extended over several metres; total lutum losses
ranged from about 18\% of the pre-degradation stock at 400 cm to more than 50\% at 60
cm, and the changes were interpreted as a divergent, initial-condition-sensitive
evolution consistent with deterministic chaos
\parencite{Montagne2013,Montagne2009}. The exported material is dominated by fine
colloids: the 0.45 $\mu$m--5000 Da fraction accounted for more than 80\% of the
suspended load, 60\% of it finer than 200 nm, and consisted mainly of 2:1
phyllosilicates with abundant small iron particles of superficial (A-horizon) origin,
so that drainage exports fine particles definitively from the profile with little
accumulation in the B horizon \parencite{Mercier2000}. These accelerated fluxes can
compress pedogenic timescales: the authors argue drainage-induced evolution over 16
years is comparable in magnitude to processes expected over much longer periods
\parencite{Montagne2010,Montagne2008}. In boreal Finland, about 70 years of improved
ditch and tile drainage was sufficient to form an argic/argillic horizon---with
illuvial clay coats maxima coinciding with the effective drainage depths---yielding
the first documented Luvisol/Cryaqualf in a boreal clay and a change judged probably
partly irreversible \parencite{YliHalla2009}.

\subsection{Macropore-controlled transport and its representation in models}

The structural findings above manifest hydrologically as rapid, connectivity-limited
preferential transport, and reproducing tile behaviour requires models that treat this
explicitly. Conservative and weakly sorbing tracers reach drains within minutes to
hours of surface application: chloride breakthrough increased abruptly by factors of
20--400 within hours, with the initial peak representing only 0.1--0.3\% of applied
mass \parencite{Villholth1998}; bromide appeared in tile effluent within 1 h and a
rainfall intensity near 2 mm h$^{-1}$ was the threshold above which concentrations
rose, yet cumulative capture was slow overall---8\% in 48 h, 27\% in 21 days, and
about 95\% only after roughly seven months \parencite{Frey2012}. Similarly, bromide
arrived at the drain as soon as drainage began even when offset from the trench, but
the bulk of transport was slow (10\% of load in 3--5 months, 50\% in 10--18 months,
90\% in 18--64 months), with travel time set by drain spacing and relative distance
from the drain \parencite{Kladivko2023}. Liquid-manure constituents (dissolved
reactive P, NH$_4^+$, Cl$^-$) reached the drain within 5 min of surface application,
and column experiments confirmed that soluble P moves through artificial macropores
with little sorption while matrix-flow columns strongly retain it---identifying
macropore flow as the primary mechanism for rapid P delivery and showing that ploughing
manure in disrupts macropores and promotes retentive matrix flow
\parencite{Stamm2002,Geohring2001}. Controlled column and laboratory experiments with
buried and surface-connected macropores quantified the geometry of these effects:
surface-connected macropores produced much faster drain response than buried ones (a
75-cm surface-connected macropore raised steady-state infiltration about fivefold),
macropore--matrix exchange was largely restricted to the profile near the drain, and
a linear total-flow-ratio relation implied an effective contributing distance of only
20--25 cm for directly connected macropores \parencite{Akay2007,Akay2008}.

These behaviours are captured only by models that represent two-domain flow, a
tile-induced two-dimensional flow field, or a direct macropore--drain connection. A
two-dimensional mobile--immobile model matched bromide breakthrough, level, and
short-term fluctuation far better than one-dimensional or single-porosity
formulations, with over 60\% of applied bromide initially transferred to immobile
regions and retained through the drainage season, indicating that both the 2-D
tile-flow field and strong mobile--immobile mass transfer control transport
\parencite{Köhne2006}. Including the macropore flow domain was essential to match
observed drainage chloride, and a one-domain simulation underpredicted tracer loss
(10\% versus the observed 21\%), while a time-variable macropore-spacing parameter
improved the chemograph fit \parencite{Villholth1998b}. In RZWQM, adding an
``express fraction'' of about 2\% reproduced the immediate breakthrough of conservative
tracer and pesticide that the unmodified model missed, improving early peaks at
multiple drained sites, though recession-limb concentrations remained poorly matched,
implying the modification only partially accounts for macropore--drain connectivity
\parencite{Fox2004,Fox2007b}. Activating a macropore component improved hourly drainage
peak amplitude but did not improve total event drainage, and predicted macropore flow
was relatively insensitive to macroporosity and pore radius \parencite{Li2023d}. A
macropore subroutine in APEX improved tile and total discharge calibration and raised
dissolved-reactive-P load performance from unsatisfactory to very good, with simulated
median annual macropore-flow and macropore-P-load contributions of 48\% and 43\%,
concentrated in winter and spring \parencite{Ford2017}. At watershed scale, tile
drainage and macropore flow were sensitive to a macropore-connectivity factor, with
fine-textured soils generating more macropore flow than coarse-textured soils, so that
soil texture explained spatial and connectivity-driven seasonal variation in flow
partitioning \parencite{Poon2021}. Biopore configuration further modulates leaching:
drain-connected biopores produced fast bromide leaching and high mass fluxes, whereas
matrix-terminating biopores reduced maximum leaching by routing some solute into the
soil \parencite{Holbak2021}. Field-scale modelling also confirms that layer-specific
macropore volumes and trench-backfill conductivity must be represented to reproduce
groundwater levels and weakly preferential tracer dynamics
\parencite{AURA1995,Bechtold2007}. Together these transport and modelling results
close the loop with the structural evidence: the small, depth-declining, drain-focused
macropore network documented morphologically is the same network that governs the
speed, magnitude, and management sensitivity of solute delivery to tile drains.

\printbibliography
\end{document}
